% !PRINTED none
% Chapter title page: no running head/page number printed on this page
% (chapter-opening convention); full body text translated below.

\begin{center}
\gesperrt{Nineteenth Chapter.}

\vspace{1em}
{\Large Calculation of the Class Invariants.}

\vspace{0.5em}
\rule{2cm}{0.4pt}

\vspace{1em}
\S\,125. \gesperrt{The Class Invariant $\gamma_2$.}
\end{center}

We have already remarked above that, besides the function $j(\omega)$,
other modular functions are also suited to the formation of class
invariants, and often give occasion to simpler results. Among these
there first presents itself the function $\gamma_2(\omega)$, which is
defined by $\sqrt[3]{j(\omega)}$. % [sic]

We have seen in \S\,71 that between the functions

\weq{1}{u = \gamma_2(\omega), \quad v = \gamma_2\left(\frac{c+\partial\omega}{a}\right),}

if $a\partial = n$ is not divisible by 3, and $c$ is assumed to be
divisible by 3, a transformation equation

\weq{2}{\Phi_n(u, v) = 0}

subsists. But the function $\Phi_n$ depends only on the two arguments
$vu^{-n}$, $u^3$, and hence in $\Phi_n(u, u) = 0$ there occurs, if
$n \equiv 1\ (\mathrm{mod}\ 3)$, only $u^3$, that is, $j(\omega)$; in
this case therefore this equation can yield no other result than the
invariant equation. It is otherwise, however, in the case

\weq{3}{n \equiv -1\ (\mathrm{mod}\ 3),}

which we shall now assume.

If one of the values is to become $v = u$, a relation must subsist

\weq{4}{\frac{\gamma+\delta\omega}{\alpha+\beta\omega} = \frac{c+\partial\omega}{a}, \quad \alpha\delta-\beta\gamma = 1,}

and at the same time [\S\,54, (15)]

\weq{5}{-\alpha\beta+\alpha\gamma+\beta\delta-\alpha^2\beta\delta \equiv 0\ (\mathrm{mod}\ 3)}

must hold. Thus in any case $\omega$ is a root of a quadratic
equation:

\weq{6}{A\omega^2+B\omega+C = 0, \quad D = B^2-4AC.}
